\section{Introduction}\label{sec:intro}

On 20 November 2026 Levent Alp\"oge released a proof of the Riemann Hypothesis
\cite{riemann1859anzahl}, and of the generalised Riemann Hypothesis for the $L$-functions of
Dirichlet \cite{dirichlet1837beweis}, formalised in Lean~4
\cite{demoura2021lean4} on top of Mathlib \cite{mathlib2020}. The release consists of the Lean
source \cite{alpoge2026lean}, \filecount{} Lean files and \linecount{} lines, and a four-page note
\cite{alpoge2026remark}
with a timeline of twenty steps. The note says that no human-readable proof has been written.

This paper does two things. It reports that the proof is correct as a formal object: we built
the archive, it compiles, its five top theorems depend on the three standard axioms and
nothing else, and the kernel accepts every theorem, and every definition that is not marked
unsafe or partial, of the import closure, replayed from an empty environment by Lean's own
checker and by a program of ours (\cref{sec:verification}). And it says
what the argument
is: an account read from the Lean source, from the final theorem down to the theorems of the
published literature on which it rests. The release does not contain one.

\subsection{The argument}
Fix a primitive Dirichlet character $\chi$. The proof has four movements.

\emph{Reduction to a one-sided bound on windows.} A windowed form of Landau's theorem reduces
the non-vanishing of $L(s,\chi)$ for $\re s>\tfrac12$ to a bound $G(x)\le x^{1/2+\eps}$ for a
smoothed sum over prime powers, required only for $x$ in long windows (\cref{sec:landau}).

\emph{An exact identity.} For an auxiliary abelian CM field with CM type $\Phi$, a weighted
count $\ASK$ of the prime powers whose Frobenius lies outside $\Phi$ equals a main term
$\MAIN$ plus a combination of such sums (\cref{sec:identity}). After averaging, an
\emph{upper} bound
\[
\ASK\ \le\ \MAIN+\Err
\]
becomes the bound on $G$, and the error has three parts of sizes $x/m$, $m$ and a rank term,
where $m$ is a parameter. At $m\approx\sqrt x$ they balance. This is the only place where
$\tfrac12$ is produced (\cref{sec:balance}).

\emph{A squeeze between two slopes.} The upper bound is proved with Hermitian lattices. For
each of many indices $u$ there is a lattice $M_u$ of algebraic modular forms on a definite
unitary group, and a map $B_u$ sending a form to the periods of its $m$-fold iterated theta
lift; $m$ is the number of lifts. At a prime $p$ inert in the CM field those periods are
divisible by about $p^{a/p+a/p^3+\cdots}$, where $a$ is the order of the Taylor coefficients
of theta functions through which they are expressed, because a CM elliptic curve has a formal
group of height two there (\cref{sec:divisibility}). So $B_u$ still maps into the lattice $P$
of periods a lattice $N_u$, obtained by rescaling a sublattice of $M_u$, whose slope exceeds
that of $M_u$ by at least the count $\ASK$, up to errors and to correction terms that cancel
later. The maps $B_u$ form a nearly tight frame, by a relative trace
formula for Bessel--Poincar\'e series (\cref{sec:frame}), and a slope inequality for such
frames bounds the slope of $P$ from below by the mean slope of the $N_u$ (\cref{sec:slope}).
From above, the slope of $P$ is bounded by $m$ times $\log(a!/\Omega^a)$, with $\Omega$ a
CM period, less a term from an auxiliary level prime; a dimension count for that prime and
the Chowla--Selberg formula for the mean of $\log\Omega$ turn this into $\MAIN$ and a
remainder that the correction terms cancel (\cref{sec:period}).

\emph{Three inputs.} The construction needs a bound on Galois cohomology in a cyclotomic tower,
which the note calls a weakening of Vandiver's conjecture (\cref{sec:cohomology}); the
existence of suitable auxiliary fields at every scale in the windows (\cref{sec:hosts}); and a
body of theorems on unitary groups, Shimura varieties and Arakelov geometry, each of which the
archive proves for itself (\cref{sec:literature}).

The shape $x/m+m$ is that of Stepanov's method \cite{stepanov1969,bombieri1974,schmidt1976equations}, where a
polynomial is made to vanish to order $m$ at the points to be counted. Here the multiplicity is
the length of a chain of theta lifts, and the degree bound is replaced by slopes of lattices in
the sense of Bost \cite{bost1996,bost2001,chambertloir2002theoremes}.

\subsection{Findings}
Beyond the outline, the reading establishes the following.
\begin{enumerate}[label=(\roman*)]
\item The statement proved is Mathlib's \lean{RiemannHypothesis}, and the statements for
  Dirichlet $L$-functions cover every character of every modulus (\cref{sec:statements}).
\item The final theorem uses \coneprojtheorems{} of the archive's \projtheorems{} theorems.
  Outside its cone are earlier forms of several steps, some with their proofs, and conditional
  statements, among them one that would give only $\re s>\tfrac34$ (\cref{sec:offpath}).
\item The central inequality is a theorem about Hermitian lattices alone, a slope inequality
  for nearly tight frames of maps, which we state in full (\cref{st:F}).
\item The period side returns the analytic main term identically; there is no discrepancy term
  (\cref{st:key}).
\item The slope argument depends on arithmetic intersection theory on a Shimura variety,
  through the covolume of lattices of integral forms; this is why resolution of singularities
  and an arithmetic Hilbert--Samuel theorem are among the results the archive formalises
  (\cref{sec:period}).
\item The cohomological bound is needed only in the weak form ``at most $\exp(x^\eps)$''. The
  archive proves more, by a route through congruences between cusp forms and Eisenstein
  classes on $\U(3)$; no statement of the main conjecture of Iwasawa theory appears among the
  results it takes from the literature (\cref{sec:cohomology}).
\item The Galois representations the proof needs are constructed in the archive from theta
  lifts between definite unitary groups, with less local--global compatibility than the
  literature provides, and the argument is arranged to need no more (\cref{sec:galois}).
\item The proof contains its own formalisations of \rowspublished{} named results of the
  literature, among them global class field theory, the Siegel--Weil formula and Hironaka's
  theorem (\cref{sec:literature}); beside these the archive proves the theorem of Ferrero and
  Washington and Kato's Euler system (\cref{sec:cohomology}).
\end{enumerate}

\subsection{Scope}
This is an exposition of Alp\"oge's proof, written from its source; the mathematics is his and
that of the model he worked with. We do not reprove the theorem, and the certificate of its
truth is the kernel's, not ours. \Cref{sec:method} states how the reading was done, and
\cref{app:audit} records where each statement of this paper was read.
